\chapter{On-Shell Scattering Amplitudes}
\label{ch:amplitudes}

% Scattering amplitudes occupy an unusual position in quantum field theory. They are among the most
% concrete objects the theory produces, feeding directly into every cross section measured at a
% collider, and yet the standard route to computing them passes through quantities that are neither
% measurable nor unique. One chooses a Lagrangian, fixes a gauge, and expands the result into Feynman
% diagrams whose individual terms depend on that gauge and on the off-shell continuation of the
% fields. Only their sum is physical. For processes with many external legs the mismatch becomes
% extreme, and thousands of diagrams routinely collapse into an answer that fits on a single line. The
% tree-level amplitude for $n$ gluons of which two carry negative helicity, discussed in
% Section~\ref{sec:BCFW}, is the classic illustration of this phenomenon.

% The on-shell program treats such simplifications as a hint rather than an accident. Rather than
% assembling amplitudes out of off-shell building blocks, one works throughout with physical,
% gauge-invariant data. The kinematics is encoded in helicity spinors, which solve the massless
% on-shell conditions identically and keep Lorentz invariance manifest at every step. The dynamics is
% constrained by general principles instead of by a specific Lagrangian, since unitarity, locality,
% and analyticity restrict the analytic structure of an amplitude so severely that the lowest-point
% amplitudes are determined by symmetry alone, while higher-point ones follow recursively from them.
% Comprehensive treatments of this material can be found in the reviews, lecture notes, and textbooks
% of Refs.~\cite{Mangano:1990by,Dixon:1996wi,Peskin:2011in,Ellis:2011cr,Dixon:2011xs,Dixon:2013uaa,%
% Elvang:2013cua,Schwartz:2014sze,Henn:2014yza,Elvang:2015rqa,Cheung:2017pzi,Travaglini:2022uwo,%
% Badger:2023eqz}.

% The interest of this perspective is not merely computational. Unitarity expresses the discontinuity
% of an amplitude through products of on-shell amplitudes with fewer loops, so that ultraviolet and
% infrared information becomes accessible without ever writing down an off-shell Green's function. The
% renormalization group equations derived in the first part of this thesis rest on that observation,
% and so do the positivity bounds obtained in the second part, where analyticity and unitarity combine
% into inequalities on the coefficients of an effective Lagrangian.

% This chapter collects the material required for those applications and fixes the conventions used
% throughout. Section~\ref{sec:spinorhelicity} develops the spinor-helicity formalism, starting from
% the decomposition of a massless momentum into helicity spinors and proceeding through the angle and
% square brackets, the identities they obey, and the wavefunctions of external states of spin up to
% two. Section~\ref{sec:Smatrix} turns to the $S$-matrix itself. There we review the properties that
% any consistent amplitude must possess, and we then examine how far those properties alone can carry
% us, fixing the three-particle amplitudes completely from helicity weights and dimensional analysis
% and reconstructing higher-point amplitudes from them by recursion.

Scattering amplitudes are the physical content of the $S$-matrix and the objects from which cross
sections are computed. Their standard evaluation, however, proceeds through quantities that are
neither measurable nor unique. One fixes a Lagrangian and a gauge, and expands the result into
Feynman diagrams whose individual terms depend on that gauge and on the off-shell continuation of
the fields, only their sum being physical. This redundancy carries a computational cost, since the
number of contributing diagrams grows rapidly with the number of external legs while the final
expression is frequently compact. The maximally-helicity-violating $n$-gluon amplitude of
Section~\ref{sec:BCFW} is the standard example.

On-shell methods bypass the intermediate redundancy by working throughout with gauge-invariant,
on-shell data. The kinematics is parametrized by helicity spinors, which solve the massless
on-shell conditions identically and keep Lorentz invariance manifest. The dynamics is constrained by
general principles instead of by a specific Lagrangian, since unitarity, locality, and analyticity
determine the singularity structure of an amplitude to the extent that the three-particle amplitudes
are fixed by helicity weights and dimensional analysis alone, and the higher-point ones are
reconstructed recursively from them. The literature on the subject is vast, and the reviews, lecture
notes, and textbooks collected in
Refs.~\cite{Mangano:1990by,Dixon:1996wi,Peskin:2011in,Ellis:2011cr,Dixon:2011xs,Dixon:2013uaa,%
Elvang:2013cua,Schwartz:2014sze,Henn:2014yza,Elvang:2015rqa,Cheung:2017pzi,Travaglini:2022uwo,%
Badger:2023eqz} constitute a very partial selection of it.

The relevance of this framework for the present work goes beyond computational efficiency. Unitarity
relates the discontinuity of an amplitude to products of on-shell amplitudes of lower loop order, so
that ultraviolet and infrared information can be extracted without introducing off-shell Green's
functions. The renormalization group equations derived in the second part of this thesis are obtained
in this way, and the theoretical bounds of the third part follow from combining analyticity with
unitarity into inequalities on the coefficients of an effective Lagrangian.

This chapter fixes the conventions used throughout and collects the results needed in what follows.
Section~\ref{sec:spinorhelicity} introduces the spinor-helicity formalism, covering the
decomposition of a massless momentum into helicity spinors, the angle and square brackets together
with the identities they satisfy, and the wavefunctions of external states of spin up to two.
Section~\ref{sec:Smatrix} deals with the $S$-matrix, reviewing first the properties that any
consistent amplitude must satisfy, then the bootstrap of the three-particle amplitudes from those
properties, and finally the recursive construction of higher-point amplitudes.

\section{Spinor-Helicity Formalism\label{sec:spinorhelicity}}

In quantum field theory, especially in gauge and gravitational theories,
the spinor-helicity formalism is a powerful tool for representing scattering amplitudes.
It is based on the idea of representing the four-momenta of the particles involved in a
scattering process in terms of spinors, which carry the fundamental representations of the
Lorentz algebra and therefore constitute the most elementary building blocks for describing the
kinematic degrees of freedom of on-shell particle states.
This section provides a brief introduction to this formalism --- mainly following Refs.~\cite{Badger:2023eqz,Henn:2014yza,Brandhuber:2022qbk} --- and establishes the conventions and definitions employed throughout the thesis.

\subsection{Helicity Spinors}

The spinor-helicity formalism exploits the fact that the proper orthochronous Lorentz group $SO^+(1,3)$ is
doubly covered by $SL(2,\mathbb C)$, i.e.,
\begin{equation}
    SO^+(1,3) \cong SL(2,\mathbb C)/\mathbb Z_2\,.
\end{equation}
Moreover, the complexified algebra splits into two commuting factors
\begin{equation}
    \mathfrak{so}(1,3)_{\mathbb C}\cong\mathfrak{sl}(2,\mathbb C)\oplus\mathfrak{sl}(2,\mathbb C)\,,
\end{equation}
so that finite-dimensional representations are labeled by a pair of half-integers $(j_-,j_+)$.

As a concrete instance of this correspondence, consider a four-vector $p^\mu = (p^0,\vec p)$, which transforms under the $(\frac{1}{2},\frac{1}{2})$ representation, and the map
\begin{equation}\label{eq:pmatrix}
    p^\mu = (p^0,\vec p) \longrightarrow p^{\dot\alpha \alpha} = p^\mu  \overline \sigma_{\mu}^{\dot\alpha \alpha} = \begin{pmatrix}
        p^0+p^3 & p^1-ip^2 \\
        p^1+ip^2 & p^0-p^3
    \end{pmatrix}\,,
\end{equation}
where $\overline \sigma_\mu^{\dot\alpha \alpha} = (\mathbb 1_{2\times 2},\vec \sigma)^{\dot\alpha \alpha}$ and $\vec \sigma = (\sigma^1,\sigma^2,\sigma^3)$ is the three-vector built from the Pauli matrices
\begin{equation}
    \sigma^1 = \begin{pmatrix}
        0& 1 \\ 1 & 0
    \end{pmatrix}\,,\qquad
    \sigma^2 = \begin{pmatrix}
        0& -i \\ i & 0
    \end{pmatrix}\,,\qquad
    \sigma^3 = \begin{pmatrix}
        1& 0 \\ 0 & -1
    \end{pmatrix}\,.
\end{equation}
Let $p^\mu$ denote the four-momentum of an on-shell particle with mass $m$. In that case, the mass-shell constraint may be interpreted as a condition on the determinant of the matrix $p^{\dot\alpha\alpha}$:\footnote{In this thesis, Lorentz indices are contracted using the mostly-minus convention for the Minkowski metric $g_{\mu\nu} = \text{diag}(+1,-1,-1,-1)$.}
\begin{equation}
    m^2 = p^\mu p_\mu = \det p^{\dot\alpha\alpha}\,.
\end{equation}
Since any $2\times 2$ matrix has at most rank two, 
$p^{\dot\alpha\alpha}$ can be generically decomposed as
\begin{equation}\label{eq:generaldecomposition}
    p^{\dot\alpha\alpha}=\tildee\lambda^{\dot\alpha}\lambda^{\alpha}+\tildee \mu^{\dot\alpha}\mu^{\alpha}\,,
\end{equation}
where $\lambda,\mu$ and $\tildee\lambda,\tildee\mu$ are commuting Weyl spinors in the $(\frac{1}{2},0)$ and $(0,\frac{1}{2})$ representations of the Lorentz group, respectively.

If we focus now on massless particles, it is clear that the general decomposition in Eq.~\eqref{eq:generaldecomposition} contains redundant variables.
In fact, for $m=0$ we have that $\det p^{\dot\alpha\alpha} = 0$, which implies that $p^{\dot\alpha\alpha}$ has a vanishing eigenvalue (the other one is $2p^0$). Thus, its rank is one and its decomposition over the spectrum reads\footnote{This fact was first used for the calculation of scattering amplitudes in Refs.~\cite{Berends:1981rb,DeCausmaecker:1981jtq,Xu:1986xb,Gunion:1985vca}.}
\begin{equation}\label{eq:spin-hel}
    p^{\dot\alpha\alpha}=2p^0\,\tildee\psi^{\dot\alpha}\psi^{\alpha} = \tildee\lambda^{\dot\alpha}\lambda^{\alpha}\,,
\end{equation}
where
$\lambda^\alpha=\sqrt{2p^0}\psi^\alpha$ and $\tildee\lambda^{\dot\alpha}=\sqrt{2p^0}\tildee\psi^{\dot\alpha}$
are the \textit{helicity spinors}.
A concrete representation of the helicity spinors is given by
\begin{align}
    \lambda^\alpha = \frac{1}{\sqrt{p^0+p^3}}\begin{pmatrix}
        p^0+p^3 \\ p^1-ip^2
    \end{pmatrix}\,, \qquad
    \tildee\lambda^{\dot\alpha} = \frac{1}{\sqrt{p^0+p^3}}\begin{pmatrix}
        p^0+p^3 \\ p^1+ip^2
    \end{pmatrix}\,.
\end{align}
If we require the four-momentum to be real, the matrix $p^{\dot\alpha\alpha}$ is Hermitian and its eigenvectors satisfy
\begin{equation}\label{eq:reality}
    (\lambda ^\alpha)^* = \pm\tildee\lambda ^{\dot\alpha}\,,
\end{equation}
where the sign is the same as that of the energy $p^0$, i.e., $p^0 >0$ for outgoing particles and $p^0<0$ for incoming ones.
Instead, if we extend the definition of the four-momentum to the complex plane, then the spinors $\lambda$ and $\tildee\lambda$ are independent.

Note that the decomposition in Eq.~\eqref{eq:spin-hel} is unique up to transformations
\begin{equation}
    \label{eq:LGscaling}
    \lambda^\alpha \longrightarrow t\,\lambda^\alpha\,,\qquad
    \tildee\lambda^{\dot\alpha}\longrightarrow  t^{-1}\,\tildee\lambda^{\dot\alpha}\,,
\end{equation}
with $t\in\mathbb C^*$ for complex momenta and $t=e^{i\phi}$ a pure phase for real ones.
This freedom is the action of the \textit{little group}, i.e., the subgroup of Lorentz transformations that leaves the four-momentum of a particle
invariant.
For a massless momentum in four dimensions, the little
group is $ISO(2)$, namely the isometry group of the Euclidean plane. Its two translation-like generators
would give rise to a continuous internal label and are represented trivially, leaving
$SO(2)\cong U(1)$, which acts by a phase (see Eq.~\eqref{eq:LGscaling}), and a
single quantum number, the helicity $h$~\cite{Wigner:1939cj,Bargmann:1948ck}.

An analogous construction holds for massive particles, where $p^{\dot\alpha\alpha}$ has rank two
and factorizes as $\tildee\lambda^{\dot\alpha I}\lambda^{\alpha}_I$~\cite{Arkani-Hamed:2017jhn}. In this case the little group is $SU(2)$ and $I=1,2$ is an index of its fundamental representation. However, throughout this thesis we focus on massless external states, as we will be interested in studying the high-energy behavior of the associated scattering amplitudes.

Spinor indices are raised and lowered with the invariant anti-symmetric spinor metric tensors, defined as
\begin{align}
    \epsilon^{\alpha\beta} = \epsilon^{\dot\alpha\dot\beta} = -\epsilon_{\alpha\beta}=-\epsilon_{\dot\alpha\dot\beta}=\begin{pmatrix}
        0 & 1 \\ -1&0
    \end{pmatrix}\,.
\end{align}
The epsilon symbols with undotted and dotted indices satisfy
\begin{align}
    \epsilon_{\alpha\beta}\epsilon^{\gamma\delta}=-\delta^\gamma_\alpha\delta^\delta_\beta+\delta_\alpha^\delta \delta_\beta^\gamma \,, \qquad
    \epsilon_{\dot\alpha\dot\beta}\epsilon^{\dot\gamma\dot\delta}=-\delta^{\dot\gamma}_{\dot\alpha}\delta^{\dot\delta}_{\dot\beta}+\delta_{\dot\alpha}^{\dot\delta} \delta_{\dot\beta}^{\dot\gamma}\,,
\end{align}
respectively,
from which it follows that 
\begin{align}
    \epsilon_{\alpha\beta}\epsilon^{\beta\gamma}=\epsilon^{\gamma\beta}\epsilon_{\beta\alpha} = \delta^\gamma_\alpha\,,
\qquad
    \epsilon_{\dot\alpha\dot\beta}\epsilon^{\dot\beta\dot\gamma}=\epsilon^{\dot\gamma\dot\beta}\epsilon_{\dot\beta\dot\alpha} = \delta^{\dot\gamma}_{\dot\alpha}\,.
\end{align}
In order to raise or lower an index of a spinor quantity, adjacent spinor indices are summed over when multiplied on the left by the appropriate epsilon symbol:
\begin{align}
    \lambda_\alpha = \epsilon_{\alpha\beta}\lambda^\beta \,,\qquad \lambda^\alpha = \epsilon^{\alpha\beta}\lambda_\beta \,,\qquad
    \tildee\lambda_{\dot\alpha} = \epsilon_{\dot\alpha\dot\beta}\tildee\lambda^{\dot\beta}\,, \qquad
    \tildee\lambda^{\dot\alpha} = \epsilon^{\dot\alpha\dot\beta}\tildee\lambda_{\dot\beta}\,.
\end{align}
Moreover, if we define $\sigma^\mu_{\alpha\dot\alpha} = \epsilon_{\alpha\beta}\epsilon_{\dot\alpha\dot\beta}\overline\sigma^{\mu\dot\beta \beta} =(\mathbb 1_{2\times 2}, \vec\sigma)_{\alpha\dot\alpha}$,
we can write the completeness relations as
\begin{align}\label{eq:compl}
    \sigma_{\mu\alpha\dot\alpha}\sigma^{\mu}_{\beta\dot\beta} = 2\epsilon_{\alpha\beta}\epsilon_{\dot\alpha\dot\beta}\,, \qquad
    \overline\sigma_{\mu}^{\dot\alpha\alpha}\overline\sigma^{\mu\dot\beta\beta} = 2\epsilon^{\alpha\beta}\epsilon^{\dot\alpha\dot\beta}\,.
\end{align}
Additionally, the identity
\begin{equation}
    \sigma^\mu_{\alpha\dot\alpha}\overline\sigma_\mu^{\dot\beta \beta}=2\delta_\alpha^\beta \delta_{\dot\alpha}^{\dot\beta}
\end{equation}
holds, and the normalization condition is given by
\begin{equation}
    \Tr(\overline\sigma_\mu \sigma_\nu)=\epsilon_{\alpha\beta} \epsilon_{\dot\alpha\dot\beta}\overline\sigma_{\mu}^{\dot\alpha\alpha} \overline\sigma_\nu^{\dot\beta\beta}=\epsilon^{\alpha\beta}\epsilon^{\dot\alpha\dot\beta}\sigma_{\mu\alpha\dot\alpha}\sigma_{\nu\beta\dot\beta}=2g_{\mu\nu}\,.
\end{equation}
Using these relations, one can demonstrate that, given the map between Lorentz four-vectors $V^\mu$ and matrices $V^{\dot\alpha\alpha}$
\begin{equation}
    V^\mu \longrightarrow V^{\dot\alpha\alpha} =  \overline\sigma^{\dot\alpha\alpha}_\mu V^\mu
\end{equation}
as the one introduced in Eq.~\eqref{eq:pmatrix}, its inverse is provided by
\begin{equation}
    V^{\dot\alpha \alpha} \longrightarrow V^\mu = \frac{1}{2}\sigma^\mu_{\alpha\dot\alpha} V^{\dot\alpha \alpha}\,.
\end{equation}
Similarly, the inverse of
\begin{equation}
    V^\mu \longrightarrow V_{\alpha\dot\alpha} =  \sigma_{\mu\alpha\dot\alpha} V^\mu
\end{equation}
is given by
\begin{equation}
    V_{\alpha\dot\alpha} \longrightarrow
    V^\mu = \frac{1}{2}\overline\sigma^{\mu\dot\alpha\alpha}V_{\alpha\dot\alpha}\,.
\end{equation}

\subsection{Angle and Square Inner Products}

In the spinor-helicity formalism, a general scattering amplitude involving massless particles is a function of the set of helicity spinors $\{\lambda_i,\tildee\lambda_i\}_{i=1}^n$, with the index $i$ running over all incoming and outgoing particles.
%Henceforth, in order to standardize the description, we will consider all particles as outgoing, i.e., $p_i^0>0$ for all $i=1,\dotsc,n$.
% Henceforth, in order to standardize the description, we will treat all particles as outgoing.
% Momentum conservation then forbids all the energies from being simultaneously positive: the legs
% that are physically incoming are those with $p_i^0<0$, and are reached by the analytic continuation
% described below.
% Amplitudes involving incoming particles are
% then obtained by crossing symmetry, which analytically continues an outgoing leg of momentum
% $p_k$ into an incoming one of momentum $-p_k$.
Henceforth, in order to standardize the description, we will treat all particles as outgoing. 
%This is a labeling convention rather than a physical restriction: 
This is not a physical restriction but just a labeling convention:
momentum conservation forbids all the energies from being simultaneously positive, and the legs that are physically incoming are simply those with $p_i^0<0$.
They are reached by crossing symmetry, which analytically continues an
outgoing leg of momentum $p_k$ into an incoming one of momentum $-p_k$.
At the level of the helicity spinors, the
flip $p_k \to p_{-k} = -p_k$ is implemented by the symmetric convention
$\lambda_k \to \lambda_{-k} = i\lambda_k$ and
$\tildee\lambda_k \to \tildee\lambda_{-k} = i\tildee\lambda_k$, which treats the two spinors on the same footing.

Starting from the helicity spinors instead of momenta, we can construct Lorentz-invariant quantities through the \textit{angle} $\agl{\cdot\,}{\cdot}$ and \textit{square} $\sqr{\cdot\,}{\cdot}$ \textit{inner products}, defined as
\begin{align}
    \langle ij\rangle= \lambda_i^\alpha \lambda_{j\alpha} = \epsilon_{\alpha\beta}\lambda_i^{\alpha}\lambda_j^{\beta} \,,\qquad
    [ij] = \tildee\lambda_{i\dot\alpha}\tildee\lambda_{j}^{\dot\alpha}=-\epsilon_{\dot\alpha\dot\beta}\tildee\lambda_i^{\dot\alpha}\tildee\lambda_j^{\dot\beta}\,,
\end{align}
with $i,j=1,\dotsc,n$, where we have introduced the north-west-to-south-east
(south-west-to-north-east) contraction rule for the undotted (dotted) Weyl indices.
Since the helicity spinors are bosonic and the epsilon symbols are anti-symmetric, the angular and square brackets are anti-symmetric as well under the exchange of their entries:
\begin{align}
    \langle ij \rangle =-\langle ji \rangle \,, \qquad
    [ij] =-[ji] \,.
\end{align}
Consequently, we have that $\langle ii \rangle = 0 = [ii]$.
The Mandelstam invariants can be written in terms of these brackets simply as
\begin{equation}
    s_{ij} = (p_i+p_j)^2 = 2p_i\cdot p_j = \lambda_i^\alpha \lambda_{j\alpha} \tildee\lambda_{j\dot\alpha}\tildee\lambda_{i}^{\dot\alpha} = \langle ij \rangle [ji]\,,
\end{equation}
which, for real momenta whose energies have the same sign, implies
\begin{equation}
    \agl{i}{j} = \sqrt{\abs{s_{ij}}} \,e^{i\phi_{ij}} = \sqr{j}{i}^*\,,
\end{equation}
with $\phi_{ij}\in\mathbb R$.

Throughout the entire thesis, to perform calculations and efficiently manipulate spinor-product expressions, we exploited the Mathematica package \textsc{S@M}~\cite{Maitre:2007jq}. 

\subsubsection{Useful Identities}

We list here some useful identities satisfied by the Lorentz-invariant inner products that we have just introduced.

\paragraph{Schouten identities.}

Any three two-component spinors must be linearly dependent, so they must satisfy
\begin{equation}\label{eq:lindep}
    c_1\lambda_1^\alpha + c_2\lambda_2^\alpha + c_3\lambda_3^\alpha=0\,,
\end{equation}
where at least one coefficient $c_i$ is different from zero.
Projecting both sides onto $\lambda_{3\alpha}$, we obtain
\begin{equation}
    c_1\agl{1}{3}+c_2 \agl{2}{3}=0\,,
\end{equation}
which leads to 
\begin{equation}\label{eq:c2}
    c_2  = \frac{\agl{3}{1}}{\agl{2}{3}}c_1\,.
\end{equation}
Similarly, the projection of Eq.~\eqref{eq:lindep} onto $\lambda_{2\alpha}$ gives
\begin{equation}
    c_1\agl{1}{2}+c_3 \agl{3}{2}=0\,,
\end{equation}
meaning that
\begin{equation}\label{eq:c3}
    c_3 = \frac{\agl{1}{2}}{\agl{2}{3}}c_1\,.
\end{equation}
Substituting Eqs.~\eqref{eq:c2} and \eqref{eq:c3} into Eq.~\eqref{eq:lindep} and dividing by $c_1\neq 0$, we obtain
\begin{equation}\label{eq:schouten}
    \lambda_1^\alpha \agl{2}{3} + \lambda_2^\alpha \agl{3}{1} + \lambda_3^\alpha \agl{1}{2} =0
\end{equation}
for any $\lambda_1,\lambda_2,\lambda_3$.
Clearly, the same steps can be repeated for the conjugate spinors to derive
\begin{equation}
    \tildee\lambda_1^{\dot\alpha} \sqr{2}{3} + \tildee\lambda_2^{\dot\alpha} \sqr{3}{1}+\tildee\lambda_3^{\dot\alpha} \sqr{1}{2} =0
\end{equation}
for any $\tildee\lambda_1,\tildee\lambda_2,\tildee\lambda_3$.

\paragraph{Momentum conservation.}

By Poincaré invariance, scattering amplitudes are non-zero only if the total four-momentum is conserved, i.e.,
\begin{equation}
    \sum_{i=1}^n p_i^\mu = 0\,,
\end{equation}
which implies
\begin{equation}\label{eq:momconservation}
    \sum_{i=1}^n \agl{a}{i}\sqr{i}{b} = 0
\end{equation}
for any $a,b$.

\paragraph{Levi-Civita contractions.}

We finally quote two identities \cite{Dixon:1996wi} that are useful for expressing the contraction of the Levi-Civita tensor with any four light-like four-vectors in terms of angle and square inner products:
\begin{align}
        \agl{a}{b}\sqr{b}{c}\agl{c}{d}\sqr{d}{a} &= \Tr( \frac{\mathbb 1 + \gamma_5}{2}\sls p_a \sls p_b \sls p_c \sls p_d)\nonumber \\ 
        &= \frac{1}{2}\left(s_{ab}s_{cd}-s_{ac}s_{bd}+s_{ad}s_{bc}-4i\epsilon_{\mu\nu\rho\sigma}p_a^\mu p_b^\nu p_c^\rho p_d^\sigma \right)\,,\label{eq:lc_identity}
\end{align}
and, similarly,
\begin{align}
        \sqr{a}{b}\agl{b}{c}\sqr{c}{d}\agl{d}{a} &= \Tr( \frac{\mathbb 1 - \gamma_5}{2}\sls p_a \sls p_b \sls p_c \sls p_d)\nonumber \\ 
        &= \frac{1}{2}\left(s_{ab}s_{cd}-s_{ac}s_{bd}+s_{ad}s_{bc}+4i\epsilon_{\mu\nu\rho\sigma}p_a^\mu p_b^\nu p_c^\rho p_d^\sigma \right)\,,
\end{align}
with $\epsilon^{0123} = -\epsilon_{0123} = 1$ and $\sls p = p_\mu \gamma^\mu$, where we have defined the Dirac matrices in the Weyl (chiral) representation as
\begin{equation}
    \gamma^\mu = \begin{pmatrix}
        \mathbb 0_{2\times 2} & \sigma^\mu\\
        \overline \sigma^{\mu} & \mathbb 0_{2\times 2}
    \end{pmatrix}\,,\qquad
    \gamma_5 = \frac{-i}{4!}\epsilon_{\mu\nu\rho\sigma}\gamma^\mu\gamma^\nu\gamma^\rho\gamma^\sigma = \begin{pmatrix}
        -\mathbb{1}_{2\times 2} & \mathbb 0_{2\times 2}\\
        \mathbb 0_{2\times 2} & \mathbb{1}_{2\times 2}
    \end{pmatrix}\,.
    \label{eq:weylgamma}
\end{equation}
The identities follow from the trace theorems
\begin{equation}
    \Tr(\gamma^\mu\gamma^\nu\gamma^\rho\gamma^\sigma)
    = 4\left(g^{\mu\nu}g^{\rho\sigma}-g^{\mu\rho}g^{\nu\sigma}+g^{\mu\sigma}g^{\nu\rho}\right)\,,
    \qquad
    \Tr(\gamma_5 \gamma^\mu\gamma^\nu\gamma^\rho\gamma^\sigma) = - 4i\epsilon^{\mu\nu\rho\sigma} \,.
\end{equation}

\subsection{Polarizations of Massless Particles\label{sec:polarizations}}

Having established that a massless momentum is encoded in a pair of helicity spinors, we now turn
to the wavefunctions carried by the external states themselves. A massless one-particle state is
specified by its momentum and by a single little-group label, the helicity $h$, and the
corresponding wavefunction is the object contracted with the amputated correlator to build a
scattering amplitude. 

For $h=0$ there is nothing to specify, while, for $h=\pm\frac{1}{2},\pm 1,
\pm 2$, the familiar covariant wavefunctions --- Dirac spinors, polarization vectors, and
polarization tensors --- carry redundant components, constrained by equations of motion and, for
the bosonic cases, defined only up to gauge transformations. Expressing them in terms of $\lambda$
and $\tildee\lambda$ removes this redundancy: each wavefunction becomes a product of helicity
spinors of definite little-group weight, the on-shell and transversality conditions reduce to the
vanishing of $\agl{p}{p}$ and $\sqr{p}{p}$, and the gauge freedom is relegated to an explicit
reference spinor that cancels from any physical quantity.

\subsubsection{Spin $\frac{1}{2}$}

The Dirac equations in momentum space for the positive- and negative-energy spinors read
\begin{align}
    (\sls p - m)\,u(p)=0\,, \qquad (\sls p + m)\,v(p)=0\,,
\end{align}
and, in the massless case, they coincide, $\sls p\, u(p) = \sls p\, v(p) = 0$. For $m=0$ chirality
and helicity agree, so that the chiral projectors $(\mathbb 1 \pm\gamma_5)/2$ produce states of
definite helicity,
\begin{align}
    u_\pm (p) = \frac{1}{2}(\mathbb 1 \pm\gamma_5)\,u(p)\,, \qquad
    v_\mp (p) = \frac{1}{2}(\mathbb 1 \pm\gamma_5)\,v(p) \,,
\end{align}
with the identification $u_\pm (p)=v_\mp (p)$, and on-shell states are labeled by
$|p,\pm \frac{1}{2}\rangle$. These polarizations are nothing but the helicity spinors themselves, as is made
manifest by the chiral, or Weyl, representation of the gamma matrices in Eq.~\eqref{eq:weylgamma},
in which $\sls p$ becomes
\begin{equation}
    \sls p = p_\mu \gamma^\mu = \begin{pmatrix}
        \mathbb 0_{2\times 2} & p_{\alpha \dot\alpha} \\ p^{\dot\alpha \alpha} & \mathbb 0_{2\times 2}
    \end{pmatrix}
    =
    \begin{pmatrix}
        \mathbb 0_{2\times 2} & \lambda_\alpha \tildee\lambda_{\dot\alpha}  \\ \tildee\lambda^{\dot\alpha}\lambda^\alpha & \mathbb 0_{2\times 2}
    \end{pmatrix}\,.
\end{equation}
Since $(\mathbb 1+\gamma_5)/2$ retains the lower two components and $(\mathbb 1-\gamma_5)/2$ the
upper ones, the two helicity solutions are
\begin{equation}
    u_+(p) = v_-(p) = \begin{pmatrix}
        \mathbb 0_{2\times 1} \\ \tildee\lambda^{\dot \alpha}
    \end{pmatrix}
    =|p]  \,,\qquad
    u_-(p) = v_+(p) = \begin{pmatrix}
        \lambda_\alpha \\ \mathbb 0_{2\times 1}
    \end{pmatrix}
    = \ket{p}\,,
\end{equation}
where we have introduced the convenient bra-ket notation with angle $\ket{\,\cdot\,}$ and square
$|\, \cdot\,]$ brackets. The massless Dirac equations,
\begin{align}
    \sls p \ket{p} = 0 \,, \qquad \sls p\, |p] = 0\,,
\end{align}
are then trivialized, as they are equivalent to $\agl{p}{p}=0=\sqr{p}{p}$.
%
For the Dirac adjoint $\overline \psi = \psi^\dagger \gamma^0$ we have
\begin{gather}
    \overline u_+(p) = \overline v_- (p)  
    = \begin{pmatrix}
        \lambda^\alpha & \mathbb 0_{1\times 2}
    \end{pmatrix}
    = \bra{p}\,,
\qquad
    \overline u_-(p) = \overline v_+ (p) 
    = \begin{pmatrix}
        \mathbb 0_{1\times 2} & \tildee\lambda_{\dot\alpha}
    \end{pmatrix}
    = [p| \,.
\end{gather}

In this notation, the block structure of $\gamma^\mu$ makes the vanishing of the chirality-violating
bilinears immediate
\begin{equation}
    \langle k| \gamma^\mu | p \rangle = 0 = [k|\gamma^\mu|p]\,,
\end{equation}
while the surviving ones reproduce the momentum
\begin{equation}
    \langle p |\gamma^\mu|p] = \lambda^\alpha \sigma^\mu_{\alpha\dot\alpha}\tildee\lambda^{\dot\alpha} = 2p^\mu\,,
    \qquad
    [p|\gamma^\mu|p\rangle = \tildee\lambda_{\dot\alpha} \overline\sigma^{\mu\dot\alpha\alpha}\lambda_\alpha = 2p^\mu\,.
\end{equation}
For an arbitrary four-vector $q^\mu$, the same structure gives
\begin{align}
    [i|\sls q|j\rangle = \tildee\lambda_{i\,\dot\alpha}\,q^{\dot\alpha \alpha}\lambda_{j\,\alpha}\,, \qquad
    \langle j |\sls q|i] = \lambda_j^\alpha\, q_{\alpha\dot\alpha}\tildee\lambda_i^{\dot\alpha}\,,
\end{align}
which implies the identity
\begin{equation}\label{eq:idgamma}
    [i|\gamma^\mu|j\rangle = \langle j |\gamma^\mu |i]\,,
\end{equation}
and, upon inserting the completeness relations of the sigma matrices in Eq.~\eqref{eq:compl}, the
Fierz identity
\begin{gather}
    [i|\gamma^\mu |j\rangle\langle l |\gamma_\mu |k] = 2\,\sqr{i}{k}\agl{l}{j}\,.
\end{gather}

The polarization degrees of freedom of external massless fermions are thus carried entirely by the
helicity spinors $\lambda$ and $\tildee\lambda$, with $|p\rangle,\langle p|$ ($|p],[ p|$) corresponding to states with helicity $h=-\frac{1}{2}$ ($h=+\frac{1}{2}$), and every fermion bilinear reduces to angle and
square brackets.

\subsubsection{Spin $1$}

Polarization vectors $\varepsilon_{\pm}^\mu$ for spin-$1$ massless gauge bosons can also be expressed in terms of helicity spinors.
In particular, positive- and negative-helicity polarization vectors can be written as
\begin{equation}\label{eq:polvecs}
    \varepsilon^\mu_+(p,\xi) = \frac{\langle \xi |\gamma^\mu|p]}{\sqrt{2}\,\agl{\xi}{p}}\,,
    \qquad 
    \varepsilon^\mu_-(p,\xi) = \frac{[\xi|\gamma^\mu|p\rangle}{\sqrt{2}\,\sqr{p}{\xi}}
    \,,
\end{equation}
respectively, or, defining $\varepsilon_{\pm}^{\dot\alpha\alpha} = \overline \sigma_\mu^{\dot\alpha\alpha}\varepsilon_{\pm}^\mu$,
\begin{equation}\label{eq:polvecsmat}
    \varepsilon_{+}^{\dot\alpha\alpha}(p,\xi) = \frac{\sqrt{2}}{\agl{\xi}{p}} \,\tildee\lambda^{\dot\alpha}\xi^\alpha\,,
    \qquad
    \varepsilon_{-}^{\dot\alpha\alpha}(p,\xi) = \frac{\sqrt{2}}{\sqr{p}{\xi}} \,\tildee\xi^{\dot\alpha}\lambda^\alpha\,,
\end{equation}
where $\xi,\tildee\xi$ are arbitrary reference spinors with the only requirement that they are not parallel to the helicity spinors $\lambda,\tildee \lambda$ associated with the momentum $p$, in such a way that $\agl{\xi}{p},\sqr{p}{\xi}\neq 0$.
Their role reflects the freedom associated with performing gauge transformations, and, as a consequence, they will not appear in any final expression for a scattering amplitude.
Indeed, if we consider an infinitesimal shift of the reference spinor $\xi \to \xi + \delta\xi$, the corresponding variation of $\varepsilon^\mu_+$ is given by
\begin{equation}
    \delta\varepsilon^\mu_+(p,\xi) = \sqrt{2}\frac{\agl{\xi\,}{\delta\xi}}{\agl{\xi}{p}^2}\,p^\mu\,,
\end{equation}
which is correctly proportional to the momentum itself.
Additionally, one can verify that the expressions in Eq.~\eqref{eq:polvecs} satisfy the properties required for the polarization vectors, namely
\begin{equation}
    (\varepsilon^\mu_{\pm}(p,\xi))^* = \varepsilon^\mu_{\mp}(p,\xi)\,,
    \qquad
    \varepsilon_\pm(p,\xi) \cdot \varepsilon_{\pm} (p,\xi) = 0\,,
    \qquad
    \varepsilon_+(p,\xi) \cdot \varepsilon_{-} (p,\xi) = -1\,,
\end{equation}
and trivialize the transversality conditions
\begin{equation}
    p\cdot \varepsilon_+(p,\xi) = \frac{1}{\sqrt{2}}\sqr{p}{p}=0\,,
    \qquad
    p\cdot \varepsilon_-(p,\xi) = \frac{1}{\sqrt{2}}\agl{p}{p}=0\,.
\end{equation}
The completeness relation
\begin{equation}
    \sum_{h=\pm} (\varepsilon^\mu_h(p,\xi))^* \varepsilon^\nu_h(p,\xi) = -g^{\mu\nu}+\frac{p^\mu \xi^\nu + p^\nu \xi^\mu}{p\cdot \xi}\,,
\end{equation}
where $\xi_\mu = \frac{1}{2}\overline\sigma_\mu^{\dot\alpha\alpha}\xi_\alpha\tildee\xi_{\dot\alpha}$,
is also satisfied.

\subsubsection{Spin $2$}
The same construction extends to massless spin-$2$ states, whose polarizations
$\varepsilon^{\mu\nu}_{\pm\pm}$ are symmetric rank-two tensors carrying helicity $h=\pm 2$.
No new ingredients are needed since they are obtained by squaring the polarization vectors of
Eq.~\eqref{eq:polvecs},
\begin{equation}\label{eq:polgrav}
    \varepsilon^{\mu\nu}_{++}(p,\xi) = \varepsilon^{\mu}_{+}(p,\xi)\,\varepsilon^{\nu}_{+}(p,\xi)\,,
    \qquad
    \varepsilon^{\mu\nu}_{--}(p,\xi) = \varepsilon^{\mu}_{-}(p,\xi)\,\varepsilon^{\nu}_{-}(p,\xi)\,,
\end{equation}
or, defining $\varepsilon_{\pm\pm}^{\dot\alpha\alpha\dot\beta\beta} = \overline \sigma_\mu^{\dot\alpha\alpha}\overline\sigma_\nu^{\dot\beta\beta}\varepsilon_{\pm\pm}^{\mu\nu}$,
\begin{equation}
    \varepsilon_{++}^{\dot\alpha\alpha\dot\beta\beta}(p,\xi)
    = \frac{2}{\agl{\xi}{p}^2}\,\tildee\lambda^{\dot\alpha}\xi^\alpha\,\tildee\lambda^{\dot\beta}\xi^\beta\,,
    \qquad
    \varepsilon_{--}^{\dot\alpha\alpha\dot\beta\beta}(p,\xi)
    = \frac{2}{\sqr{p}{\xi}^2}\,\tildee\xi^{\dot\alpha}\lambda^\alpha\,\tildee\xi^{\dot\beta}\lambda^\beta\,,
\end{equation}
with the same reference spinors $\xi,\tildee\xi$ and the same requirement
$\agl{\xi}{p},\sqr{p}{\xi}\neq 0$.
The symmetry in the Lorentz indices $\mu\leftrightarrow\nu$ is manifest, while the remaining defining properties are
inherited from the spin-$1$ case,
\begin{equation}
    (\varepsilon^{\mu\nu}_{\pm\pm}(p,\xi))^* = \varepsilon^{\mu\nu}_{\mp\mp}(p,\xi)\,,
    \qquad
    \varepsilon_{++}(p,\xi) \cdot \varepsilon_{--}(p,\xi) = 1\,,
\end{equation}
and both transversality and tracelessness are trivialized,
\begin{equation}
    p_\mu\varepsilon^{\mu\nu}_{\pm\pm}(p,\xi) = (p\cdot\varepsilon_\pm(p,\xi))\varepsilon^\nu_\pm(p,\xi) = 0\,,
    \qquad
    g_{\mu\nu}\varepsilon^{\mu\nu}_{\pm\pm}(p,\xi) = \varepsilon_\pm(p,\xi)\cdot\varepsilon_\pm(p,\xi) = 0\,,
\end{equation}
so that no traceless projection has to be imposed by hand.
The reference spinors drop out of physical quantities exactly as before: under
$\xi \to \xi + \delta\xi$ one has $\delta\varepsilon^\mu_+ \propto p^\mu$, and therefore
\begin{equation}
   \delta\varepsilon^{\mu\nu}_{++}(p,\xi) = \sqrt{2}\frac{\agl{\xi\,}{\delta\xi}}{\agl{\xi}{p}^2}\left(p^\mu \varepsilon_+^\nu + p^\nu \varepsilon_+^\mu \right)\,,
\end{equation}
which is a linearized diffeomorphism, the spin-$2$ counterpart of the gauge transformation
encountered for spin $1$, and it leaves scattering amplitudes invariant.

\section{\texorpdfstring{$S$-Matrix}{S-Matrix}\label{sec:Smatrix}}

% The \textit{scattering operator} $S$, or \textit{$S$-matrix}, evolves an asymptotic state defined at $t=-\infty$, $\ket{\psi}_{\text{in}}$, to a different asymptotic state defined at $t=+\infty$, $\ket{\chi}_{\text{out}}$:
% \begin{equation}
%     \ket{\chi}_{\text{out}} = S \ket{\psi}_{\text{in}}\,.
% \end{equation}
% The associated matrix elements are what we refer to as \textit{scattering amplitudes}:
% \begin{equation}
%     S(\psi \to \chi) = \mel{\chi}{S}{\psi}\,.
% \end{equation}

The \textit{scattering operator} $S$, or \textit{$S$-matrix}, connects the asymptotic states
prepared at $t=-\infty$ with those observed at $t=+\infty$. These form two complete bases of the
same Hilbert space, and $S$ is the unitary operator relating them, so that the probability amplitude
for an initial configuration $\ket{\psi}$ to be observed as a final configuration $\ket{\chi}$ is
\begin{equation}
    S(\psi \to \chi) = \mel{\chi}{S}{\psi}\,.
\end{equation}
These matrix elements are what we refer to as \textit{scattering amplitudes}.

\subsection{Properties and Symmetries of Scattering Amplitudes}

Before computing any amplitude, it is worth asking what an amplitude is allowed to be. The
$S$-matrix is not an arbitrary function of the external kinematics: a small number of physical
principles --- conservation of probability, the pointlike nature of interactions, causality, and
the symmetries of the vacuum --- constrain it so severely that in favorable cases they determine it
outright, as we will see for three-particle amplitudes in the next subsection. 
It is convenient to
separate the trivial forward propagation from the genuine interaction by writing
\begin{equation}\label{eq:Tdef}
    S = \mathbb 1 + iT\,,
\end{equation}
and to strip from $T$ the overall momentum-conserving delta function,
\begin{equation}\label{eq:Mdef}
    \mel{\chi}{T}{\psi} = (2\pi)^4\,\delta^{(4)}\!\left(\sum_{i\in\psi} p_i - \sum_{i\in\chi} p_i \right)\mathcal M(\psi\to\chi)\,,
\end{equation}
so that $\mathcal M$ is the object of interest. We review below the properties that $\mathcal M$
must satisfy.

\subsubsection{Unitarity\label{sec:unitarity}}

A defining property of the scattering operator is \textit{unitarity}, the operator statement of
probability conservation:
\begin{equation}\label{eq:unitarity}
    S^\dagger S = SS^\dagger = \mathbb 1\,.
\end{equation}
Inserting Eq.~\eqref{eq:Tdef} into Eq.~\eqref{eq:unitarity}, the linear terms cancel against each
other only up to a quadratic remainder,
\begin{equation}
    -i\left(T - T^\dagger\right) = T^\dagger T\,,
\end{equation}
which, once a complete set of intermediate states $\ket{X}$ is inserted and the delta function of
Eq.~\eqref{eq:Mdef} is stripped, reads
\begin{equation}\label{eq:cutting}
    -i\left(\mathcal M(\psi\to\chi) - \mathcal M(\chi\to\psi)^*\right)
    = \sum_X \int \text{dLIPS}_X\,\mathcal M(\chi\to X)^*\,\mathcal M(\psi\to X)\,,
\end{equation}
where the sum runs over all intermediate states allowed by the quantum numbers and by the available
energy, and $\text{dLIPS}_X$ denotes the Lorentz-invariant phase-space measure. Setting
$\chi=\psi$ gives the optical theorem
\begin{equation}\label{eq:optical}
    2\Im \mathcal M(\psi\to\psi) = \sum_X \int \text{dLIPS}_X\,\abs{\mathcal M(\psi\to X)}^2\,,
\end{equation}
whose right-hand side is manifestly non-negative.

The underlying physical interpretation is straightforward: whatever probability is removed from the forward direction must
reappear in some final state, so the absorptive part of a forward amplitude measures the total
probability of scattering into anything at all. 
%The consequence that will matter most in this thesis is, however, structural rather than numerical. 
Eq.~\eqref{eq:cutting} expresses the
\textit{imaginary} part of an amplitude entirely in terms of \textit{on-shell} amplitudes with fewer
legs or fewer loops. Indeed, the intermediate states $\ket{X}$ are physical states integrated over their
phase space, and no off-shell information whatsoever enters. Thus, unitarity turns the computation of
a complicated object into the gluing of simpler ones, 
and the on-shell methods employed in the following chapters exploit exactly this mechanism.

\subsubsection{Locality}

% Interactions among fields take place at a point in spacetime. The imprint of this fact on the
% $S$-matrix is a statement about \textit{where amplitudes are allowed to be singular}: the only
% singularities of a tree-level amplitude are simple poles occurring when an internal particle goes on
% shell. Denoting by $P=\sum_{i\in I}p_i$ the total momentum flowing through an internal line, with
% $I$ a subset of the external legs, one has
% \begin{equation}\label{eq:factorization}
%     \mathcal M_n \;\xrightarrow[P^2\to 0]{}\;
%     \sum_{h}\,\mathcal M_{L}\!\left(\{i\in I\},\,P^{h}\right)\,
%     \frac{1}{P^2}\,
%     \mathcal M_{R}\!\left(-P^{-h},\,\{i\notin I\}\right)\,,
% \end{equation}
% where the sum runs over the helicity states of the exchanged particle. The residue at the pole is the product of two lower-point on-shell amplitudes, glued
% along the internal line.

% Intuitively, a pole at $P^2=0$ is the signature of an intermediate particle travelling a
% macroscopically long distance before interacting again, so that the process factorizes into two
% causally separated sub-processes. Anything else --- a double pole, a pole at a value of $P^2$ that
% does not correspond to a physical state, or an essential singularity --- would describe a
% non-local interaction, in the same sense in which an operator with inverse powers of derivatives
% does. Equations~\eqref{eq:cutting} and \eqref{eq:factorization} act as the two faces of the same
% constraint, one on the discontinuities and one on the poles, and together they are strong enough to
% reconstruct amplitudes recursively, as discussed in Section~\ref{sec:BCFW}.

Locality expresses the fact that interactions take place at well-defined points of spacetime. Its
imprint on the $S$-matrix is a constraint on the position of the singularities, since an amplitude
can become singular only when an intermediate state goes on shell and propagates over macroscopic
distances. The exchange of a single particle of mass $m$ produces a simple pole at $P^2=m^2$, while
a multi-particle intermediate state produces a branch point at the corresponding threshold. The
analytic structure of the $S$-matrix therefore mirrors the spectrum of the theory.

Splitting the $n$ external legs into a subset $I$ and its complement, and denoting by
$P_I=\sum_{i\in I}p_i$ the momentum exchanged between the two halves, the residue at a
single-particle pole factorizes into on-shell amplitudes of lower multiplicity
\begin{equation}\label{eq:factorization}
    \lim_{P_I^2\to m^2} (P_I^2-m^2)\,\mathcal M_n
    = \sum_{h}\mathcal M_{\abs{I}+1}\Big(\{i\}_{i\in I},(-P_I)^{-h}\Big)\,
      \mathcal M_{n-\abs{I}+1}\Big(P_I^{h},\{j\}_{j\notin I}\Big)\,,
\end{equation}
where the sum runs over the helicity states of the exchanged particle.

Locality and unitarity are complementary in this respect. The former dictates where an amplitude is
singular, the latter dictates how it behaves there. Eq.~\eqref{eq:factorization} fixes the
residue at a single-particle pole through a plain product of lower-point amplitudes, while the
discontinuity across a multi-particle cut requires in addition an integration over the
phase space of the intermediate state, in the form already encountered in
Eq.~\eqref{eq:cutting}. Both statements support the recursive construction of higher-point
amplitudes out of simpler ones developed in Section~\ref{sec:BCFW}, and both are the
momentum-space counterpart of the cluster decomposition principle, according to which distant sub-processes remain correlated only
through the exchange of on-shell particles.

\subsubsection{Analyticity\label{sec:analyticity}}

Locality and unitarity acquire their full strength once the kinematic invariants are allowed to
become complex. Causality, the requirement that no signal propagate outside the light
cone, implies that amplitudes are boundary values of functions analytic in the complexified
Mandelstam invariants, with singularities confined to the positions dictated by the physical
spectrum. Poles sit at the masses of stable one-particle states and branch cuts open at
multi-particle thresholds, in accordance with the discussion of locality above. The physical
amplitude is obtained by approaching the real axis from above, an instruction encoded in the
Feynman $i\varepsilon$ prescription of the propagators.

Two consequences deserve to be singled out. The first one is crossing symmetry. A single analytic
function describes several physically distinct processes at once, and one passes from one to another
by continuing an external leg between the initial and the final state, in the sense already employed
in Section~\ref{sec:spinorhelicity} when adopting the all-outgoing convention.

The second consequence is quantitative. Analyticity allows Cauchy's theorem to be applied to a small
contour encircling the low-energy region of the complexified center-of-mass energy, which can then
be deformed onto the cuts running along the real axis, provided the amplitude is bounded well enough
at large energies for the arcs at infinity to be discarded. The outcome is a family of sum rules, in
which the coefficients of the low-energy expansion of the amplitude, computable within an effective
field theory in terms of its Wilson coefficients, are expressed as integrals of the discontinuity
over the entire high-energy region. Unitarity then enters through Eq.~\eqref{eq:optical}, which
forces that discontinuity to be non-negative, so the sum rules turn into inequalities. The low-energy
couplings of any effective theory admitting a causal and unitary ultraviolet completion are thereby
constrained to lie inside a convex cone~\cite{Arkani-Hamed:2020blm,Adams:2006sv}, and analyticity
converts information about physics that one cannot resolve into statements testable at energies one
can reach.

\paragraph{Dispersion relations and positivity bounds.}
For the derivation we follow Refs.~\cite{Adams:2006sv,Bellazzini:2015cra}. Consider the elastic $2\to2$ amplitude
$\mathcal M(s,t)$ of massless states, with
\begin{equation}
    s=(p_1+p_2)^2\,,\qquad t=(p_1+p_3)^2\,,\qquad u=(p_1+p_4)^2\,.
\end{equation}
Crossing in the $t$-channel exchanges $s\leftrightarrow u$ at fixed $t$, so that
\begin{equation}\label{eq:crossingrelation}
    \mathcal M(s,t) = \mathcal M(-s-t,t)\,,
\end{equation}
which in the forward limit reduces to $\mathcal M(s,0)=\mathcal M(-s,0)$. The two operations are
compatible only for suitable external states, with particle $3$ the crossed image of particle $1$
and particle $4$ that of particle $2$, and with the polarizations held fixed relative to the momenta.

Below the first threshold of the ultraviolet completion the forward amplitude is regular and admits
the low-energy expansion
\begin{equation}\label{eq:lowenergyexp}
    \mathcal M(s,t\to0) = \sum_{n\ge 0} c_n\,s^n\,,
\end{equation}
whose coefficients are computable within the effective theory in terms of its Wilson coefficients.
They are extracted by a contour integral around the origin
\begin{equation}\label{eq:smallcontour}
    c_n = \frac{1}{2\pi i}\oint_{\mathcal C}\frac{\text{d}s}{s^{n+1}}\,\mathcal M(s,t\to0)\,,
\end{equation}
with $\mathcal C$ a small circle enclosing no singularity other than the pole at $s=0$.

\begin{figure}[tb]
\centering
\begin{tikzpicture}[scale=1.0, >=Stealth,
  cutline/.style={red!70!black, very thick, decorate,
                  decoration={zigzag, segment length=4.5pt, amplitude=1.6pt}},
  cont/.style={thick, blue!55!black},
  arr/.style={postaction={decorate},
              decoration={markings, mark=at position #1 with {\arrow{>}}}}]

  % axes
  \draw[->, gray!70] (-5.6,0) -- (5.6,0) node[right, black] {$\operatorname{Re}s$};
  \draw[->, gray!70] (0,-4.6) -- (0,4.6) node[above, black] {$\operatorname{Im}s$};

  % branch cuts
  \draw[cutline] (1.3,0) -- (5.3,0);
  \draw[cutline] (-1.3,0) -- (-5.3,0);
  \fill (1.3,0) circle (1.8pt);  \node[above=2.5pt] at (1.3,0) {$s_0$};
  \fill (-1.3,0) circle (1.8pt); \node[below=1.5pt]  at (-1.3,0) {$-s_0$};

  % small contour around the origin
  \draw[cont, arr=0.25, arr=0.75] (0.62,0) arc (0:360:0.62);
  \node[cont] at (0.62,0.62) {$\mathcal C$};

    % deformed contour: pieces hugging the two cuts
  \draw[cont, arr=0.30, arr=0.85]
        (3.9,-0.18) -- (1.3,-0.18) arc (270:90:0.18) -- (3.9,0.18);
  \draw[cont, arr=0.30, arr=0.85]
        (-3.9,0.18) -- (-1.3,0.18) arc (90:-90:0.18) -- (-3.9,-0.18);

  % deformed contour: arcs at large |s|
  \draw[cont, dashed, arr=0.5] (3.9,0.18)  arc (2.65:177.35:3.904);
  \draw[cont, dashed, arr=0.5] (-3.9,-0.18) arc (182.65:357.35:3.904);
  \node[cont] at (3.05,3.05) {$\mathcal C'$};
\end{tikzpicture}
\caption{Analytic structure of the crossing-symmetric forward amplitude $\mathcal M(s,0)$ in the complex $s$ plane. The
branch cuts start at the first threshold $s_0$ of the ultraviolet completion and at its crossed image
$-s_0$. The small contour $\mathcal C$ of Eq.~\eqref{eq:smallcontour} is deformed into $\mathcal C'$,
which wraps the two cuts and closes on arcs at large $\abs{s}$. The arcs are discarded under the
boundedness assumption discussed in the text, leaving the dispersion relation
of Eq.~\eqref{eq:dispersion}.}
\label{fig:contour}
\end{figure}

Unitarity generates a branch cut along the real axis from the first threshold $s_0$ to $+\infty$
and, by crossing, a mirror cut from $-s_0$ to $-\infty$. By Cauchy's theorem, $\mathcal C$ can be deformed outwards without crossing any singularity into the
contour $\mathcal C'$ of Figure~\ref{fig:contour}, which runs along both sides of the two cuts and
closes on a circle of large radius. The
arcs can be discarded provided $\abs{\mathcal M(s,0)}$ grows more slowly than $\abs{s}^{n}$, which
the Froissart bound guarantees for $n\ge2$ \cite{Froissart:1961ux}.

What survives is the integral of the discontinuity across the cuts. The Schwarz reflection
principle, $\mathcal M(s^{*},0)=\mathcal M(s,0)^{*}$, gives
\begin{equation}
    \operatorname{Disc}\mathcal M(s,0) = \mathcal M(s+i\varepsilon,0)-\mathcal M(s-i\varepsilon,0)
    =2i\,\Im\mathcal M(s,0)\,,
\end{equation}
while Eq.~\eqref{eq:crossingrelation} in the forward limit maps the left cut onto the right one.
Collecting the two contributions, we obtain
\begin{equation}\label{eq:dispersion}
    c_n = \frac{1+(-1)^n}{\pi}\int_{s_0}^{\infty}\frac{\text{d}s}{s^{n+1}}\,\Im\mathcal M(s,0)\,.
\end{equation}
The odd coefficients vanish, as required by crossing symmetry, and the even ones are fixed by
the absorptive part of the amplitude integrated over the entire ultraviolet region.

Finally, unitarity closes the argument: the optical theorem in Eq.~\eqref{eq:optical} relates the forward
absorptive part to the total cross section through $\Im\mathcal M(s,0)=s\,\sigma(s)$, and $\sigma(s)$
is always non-negative. Therefore
\begin{equation}\label{eq:positivity}
    c_n = \frac{2}{\pi}\int_{s_0}^{\infty}\text{d}s\;\frac{\sigma(s)}{s^{n}} \ge 0
\end{equation}
for $n\ge 2$ and even.
A low-energy coefficient associated with an effective operator of dimension $8,12,16,\dotsc$ is thus bounded by an integral over a physical region one cannot resolve, and its
sign is fixed without any knowledge of the ultraviolet completion beyond causality and unitarity.

The construction extends to superpositions of states. Given two one-particle superpositions
\begin{equation}
    \ket{u}=\sum_i u_i\ket{i}\,,\qquad \ket{v}=\sum_j v_j\ket{j}\,,
\end{equation}
over the species and polarizations available at low energy, the corresponding elastic amplitude reads
\begin{equation}\label{eq:superposition}
    \mathcal M(s,t) = \mathcal M(uv\to uv)
    = \sum_{i,j,k,l} u_i^{*}v_j^{*}u_kv_l\,\mathcal M(kl\to ij)\,.
\end{equation}
The
argument above applies to it unchanged, and Eq.~\eqref{eq:positivity} becomes
\begin{equation}\label{eq:cone}
    \left.\dv[n]{}{s}\,\mathcal M(s,0)\right|_{s=0} \ge 0
\end{equation}
for $n\ge 2$ and even, and
for every admissible choice of the coefficients $u_i,v_j$. The
left-hand side is quartic in the superposition coefficients and linear in the Wilson coefficients of
the effective theory, so each choice supplies one linear inequality among them, and the whole family
carves out the convex cone anticipated above. This is the mechanism behind the positivity bounds
derived in the third part of this thesis.

\subsubsection{Poincaré Transformations}

Invariance of the vacuum under spacetime translations forces the total four-momentum to be
conserved. This conservation is already accounted for by the delta function factored out in
Eq.~\eqref{eq:Mdef}, and at the level of the helicity spinors it takes the form of
Eq.~\eqref{eq:momconservation}. Invariance under
Lorentz transformations forces $\mathcal M$ to depend on the external data only through Lorentz
invariant quantities, which, for massless particles, are represented by the angle and square brackets defined above.

In fact, an infinitesimal Lorentz transformation acts as
\begin{align}\label{eq:lorentztransf}
    \delta \lambda^\alpha = \omega^{\alpha\beta}\lambda_\beta \,, \qquad
    \delta\tildee\lambda_{\dot\alpha} = \tildee\omega_{\dot\alpha\dot\beta}\tildee\lambda^{\dot\beta}\,,
\end{align}
with $\omega^{\alpha\beta}$ and $\tildee\omega_{\dot\alpha\dot\beta}$ symmetric. Their symmetry property is equivalent to the tracelessness of the corresponding
$\mathfrak{sl}(2,\mathbb C)$ generators, since the anti-symmetric part of a rank-two spinor is
proportional to the epsilon symbol. The two transformations are generated by
$\tfrac{1}{2}\omega^{\alpha\beta}\mathbb M_{\alpha\beta}$ and
$\tfrac{1}{2}\tildee\omega_{\dot\alpha\dot\beta}\tildee{ \mathbb M}^{\dot\alpha\dot\beta}$, where the
\textit{angular-momentum operators} are defined as
\begin{gather}
    \mathbb M_{\alpha\beta} = \sum_{i=1}^n \left(\lambda_{i\alpha}\frac{\partial}{\partial \lambda_i^\beta}
    + \lambda_{i\beta}\frac{\partial}{\partial \lambda_i^\alpha}\right)\,,\qquad
    \tildee{\mathbb M}^{\dot\alpha\dot\beta} = \sum_{i=1}^n \left(\tildee\lambda_i^{\dot\alpha}\frac{\partial}{\partial \tildee\lambda_{i\dot\beta}}
    + \tildee\lambda_i^{\dot\beta}\frac{\partial}{\partial \tildee\lambda_{i\dot\alpha}}\right)\,.
\end{gather}
The brackets are annihilated by both:
\begin{equation}
    \delta\agl{i}{j} = \omega^{\alpha\beta}\Big(\lambda_{i\beta}\lambda_{j\alpha}
    - \lambda_{i\alpha}\lambda_{j\beta}\Big) = 0\,,
    \qquad
    \delta\sqr{i}{j} = \tildee\omega_{\dot\alpha\dot\beta}\Big(\tildee\lambda_i^{\dot\beta}\tildee\lambda_j^{\dot\alpha}
    - \tildee\lambda_i^{\dot\alpha}\tildee\lambda_j^{\dot\beta}\Big) = 0\,,
\end{equation}
where in each case the second term reproduces the first once the summed indices are relabeled, and
the two cancel by the symmetry of $\omega$ and $\tildee\omega$. 
Lorentz invariance of the
$S$-matrix therefore amounts to the statement that the amplitude depends on the external kinematics
only through the brackets,
\begin{equation}
    \mathcal M = \mathcal M\big(\{\agl{i}{j},\sqr{i}{j}\}\big)\,.
\end{equation}

\subsubsection{Little-Group Transformations}

Lorentz invariance leaves an amplitude a function of the angle and square brackets alone, and yet
the brackets themselves are not invariant under the rescaling of Eq.~\eqref{eq:LGscaling}.
Amplitudes are then
accordingly required to be covariant, carrying a definite weight on every external leg.

The dependence of $\mathcal M$ on $\lambda_i$ and $\tildee\lambda_i$ enters through the
wavefunctions of the external states, so the amplitude is homogeneous of definite degree in the
spinors of each leg. At tree level, it is a rational function of the brackets. Beyond tree level, the
loop integrations produce transcendental functions of the Mandelstam invariants, which are inert
under the little group, so the degree is still carried entirely by the rational spinor
prefactors. These degrees are counted by the \textit{helicity operator}
\begin{equation}\label{eq:helicityop}
    \mathbb H_i = -\frac{1}{2}\lambda_i^\alpha \frac{\partial}{\partial \lambda_i^\alpha}
    + \frac{1}{2}\tildee\lambda_{i\dot\alpha}\frac{\partial}{\partial \tildee\lambda_{i\dot\alpha}}\,,
\end{equation}
%
As an illustration, the polarization matrices of Eq.~\eqref{eq:polvecsmat} are eigenvectors with the expected eigenvalues,
\begin{equation}
    \mathbb H\,\varepsilon_{+}^{\dot\alpha\alpha}(p,\xi) = +\varepsilon_{+}^{\dot\alpha\alpha}(p,\xi)\,,
    \qquad
    \mathbb H\,\varepsilon_{-}^{\dot\alpha\alpha}(p,\xi) = -\,\varepsilon_{-}^{\dot\alpha\alpha}(p,\xi)\,.
\end{equation}

Requiring the particle $i$ to have helicity $h_i$ therefore amounts to the homogeneity condition
\begin{equation}
    \mathbb H_i\,\mathcal M\big(1^{h_1},\dotsc,n^{h_n}\big) = h_i\,\mathcal M\big(1^{h_1},\dotsc,n^{h_n}\big)
\end{equation}
for every $i=1,\dotsc,n$, or, equivalently,
\begin{equation}\label{eq:LGweight}
    \mathcal M\big(\dotsc,t_i\lambda_i,t_i^{-1}\tildee\lambda_i,\dotsc\big)
    = t_i^{-2h_i}\,\mathcal M\big(\dotsc,\lambda_i,\tildee\lambda_i,\dotsc\big)\,.
\end{equation}


\subsubsection{Scale Transformations}

With the relativistic normalization of the asymptotic states
\begin{equation}
    \braket{p}{p'} = 2E_{\vec p}\,(2\pi)^3\,\delta^{(3)}(\vec p - \vec p\,')
\end{equation}
each one-particle state carries mass dimension $-1$, and Eq.~\eqref{eq:Mdef} then fixes\footnote{In a dimension-$d$ spacetime, a generic $n$-point scattering amplitude has mass dimension $[\mathcal M_n] = d+n-dn/2$.}
\begin{equation}\label{eq:massdim}
    [\mathcal M_n] = 4-n\,.
\end{equation}

Eq.~\eqref{eq:massdim} is exact, being fixed by the normalization of the states alone. What
holds only at tree level is that the amplitude is a homogeneous function of the external momenta.
In dimensional regularization, the loop integrations introduce the renormalization scale $\mu$, and
the amplitude acquires a dependence on the dimensionless ratios of the kinematic invariants to
$\mu^2$. Being a physical quantity, the amplitude cannot depend on $\mu$, so this explicit
dependence has to cancel against the implicit one carried by the renormalized couplings. The
cancellation is what forces the couplings to run, according to their anomalous dimension, and its systematic exploitation is the subject of
the second part of the thesis.


\subsection{Bootstrap of Three-Point Amplitudes}
\label{sec:3pointamplitudes}

Following Ref.~\cite{Benincasa:2007xk}, we show here how the constraints that arise from locality, Poincaré invariance, and dimensional analysis are enough to fix the form of three-particle scattering amplitudes up to an overall coupling constant.

We can start by noting that momentum conservation
\begin{equation}
    p_1+p_2+p_3=0
\end{equation}
implies, for massless on-shell particles,
\begin{equation}
    p_1^2=\agl{2}{3}\sqr{3}{2}=0\,,\qquad
    p_2^2=\agl{1}{3}\sqr{3}{1}=0\,,\qquad
    p_3^2=\agl{1}{2}\sqr{2}{1}=0\,.
\end{equation}
Thus, we deduce that if $\agl{1}{2}\neq 0$, then $\sqr{1}{2}=0$. Moreover, applying Eq.~\eqref{eq:momconservation}, we find that $\agl{1}{2}\sqr{2}{3} = -\agl{1}{1}\sqr{1}{3}-\agl{1}{3}\sqr{3}{3} = 0$, therefore $\sqr{2}{3}=0$.
By repeating this procedure cyclically, we conclude that $\sqr{1}{2}=\sqr{2}{3}=\sqr{3}{1}=0$. 
By contrast, if we assume $\sqr{1}{2}\neq 0$, we would find that $\agl{1}{2}=\agl{2}{3}=\agl{3}{1}=0$.

To summarize, three-particle amplitudes have support on two different kinematic configurations:
\begin{itemize}
    \item[($\mathsf H$)] The \textit{holomorphic} configuration, corresponding to the case in which all square brackets are zero,
    \begin{equation}
        \sqr{1}{2}=\sqr{2}{3}=\sqr{3}{1}=0\,,
    \end{equation}
    implying $\tildee\lambda_1\propto \tildee\lambda_2 \propto \tildee\lambda_3$;
    \item[($\mathsf A$)] The \textit{anti-holomorphic} configuration, corresponding to the case in which all angle brackets are zero,
    \begin{equation}
        \agl{1}{2}=\agl{2}{3}=\agl{3}{1}=0\,,
    \end{equation}
    implying $\lambda_1\propto \lambda_2 \propto \lambda_3$.
\end{itemize}
In both configurations, the four-momenta need to be complexified, so that $\lambda_i$ and $\tildee \lambda_i$ can be treated as independent variables. For real momenta the two configurations collapse onto each other, and any amplitude with
non-trivial bracket dependence vanishes.

We can consider the following ansatz for the three-point amplitude of particles with helicities $h_i$:
\begin{equation}
    \mathcal M (1^{h_1},2^{h_2},3^{h_3}) = \mathcal M_{\mathsf H}(\agl{1}{2},\agl{2}{3},\agl{3}{1}) + \mathcal M_{\mathsf A}(\sqr{1}{2},\sqr{2}{3},\sqr{3}{1})\,,
\end{equation}
where we have defined the holomorphic and anti-holomorphic parts as
\begin{equation}
    \mathcal M_{\mathsf H}(\agl{1}{2},\agl{2}{3},\agl{3}{1}) = g_{\mathsf H} \agl{1}{2}^a \agl{2}{3}^b \agl{3}{1}^c \,,\quad
    \mathcal M_{\mathsf A}(\sqr{1}{2},\sqr{2}{3},\sqr{3}{1}) = g_{\mathsf A} \sqr{1}{2}^{-a} \sqr{2}{3}^{-b} \sqr{3}{1}^{-c}\,,
\end{equation}
respectively, and $g_{\mathsf H},g_{\mathsf A}$ are the corresponding coupling constants.
The three exponents $a,b,c$ can be fixed by imposing the homogeneity condition for each particle helicity, i.e.,
\begin{equation}
    -2h_1 = a+c\,,\qquad -2h_2=a+b\,,\qquad -2h_3=b+c\,,
\end{equation}
whose solution is
\begin{equation}
    a = h_3-h_1-h_2\,,
    \qquad
     b = h_1-h_2-h_3\,,
     \qquad
      c = h_2-h_3-h_1\,,
\end{equation}
and, in particular, $a,b,c\in\mathbb Z$ since fermions always appear in pairs.
Moreover, by requiring the amplitude to have the correct mass dimension, $[\mathcal M]=1$, we obtain
\begin{equation}
    [g_{\mathsf H}] = 1+h\,,\qquad [g_{\mathsf A}] = 1-h\,,
\end{equation}
where $h=h_1+h_2+h_3$.

The final step consists in requiring the amplitude to be well behaved in the limit of real momenta.
As the kinematic configurations $\mathsf H$ and $\mathsf A$ degenerate when the momenta are taken to
be real, all brackets vanish simultaneously, and $\mathcal M$ must remain finite there. Since the
total bracket power of the two terms is $-h$ and $h$
for the holomorphic and anti-holomorphic parts, respectively, the requirement selects one of the two
branches according to the sign of the total helicity:
\begin{itemize}
    \item If $h<0$, the anti-holomorphic term diverges and we must set $g_{\mathsf A}=0$, so that
    \begin{equation}\label{eq:3parthol}
        \mathcal M (1^{h_1},2^{h_2},3^{h_3}) = g_{\mathsf H}\, \agl{1}{2}^{h_3-h_1-h_2}\, \agl{2}{3}^{h_1 - h_2 - h_3}\, \agl{3}{1}^{h_2 - h_3 - h_1}\,,
    \end{equation}
    with $[g_{\mathsf H}] = 1+h<1$;
    \item If $h>0$, the holomorphic term diverges and we must set $g_{\mathsf H}=0$, so that
    \begin{equation}\label{eq:3partahol}
        \mathcal M (1^{h_1},2^{h_2},3^{h_3}) = g_{\mathsf A}\, \sqr{1}{2}^{h_1+h_2-h_3}\, \sqr{2}{3}^{h_2+h_3-h_1}\, \sqr{3}{1}^{h_3+h_1-h_2} \,,
    \end{equation}
    with $[g_{\mathsf A}] = 1-h<1$;
    \item If $h=0$, both terms are of degree zero in the brackets and the argument above does not
    discriminate between them. Both couplings are then allowed and have $[g_{\mathsf H}]=[g_{\mathsf A}]=1$, so that the
    corresponding operator has dimension $[\mathcal O]=3$ and contains three fields. The only such
    local operator in four dimensions is $\phi^3$, since two fermions already saturate dimension
    three by themselves and any field strength or derivative exceeds it. Locality therefore forces
    $h_1=h_2=h_3=0$ and $a=b=c=0$, leaving the constant amplitude of the cubic scalar coupling as
    the unique marginal possibility.
\end{itemize}

In all three cases, the surviving coupling has mass dimension not larger than one. This feature is
not an accident but rather reflects the underlying principle of locality.
Had we kept the divergent branch instead, its coupling
would have carried $[g]>1$, and would therefore have descended from an operator $\mathcal O$ of mass
dimension $[\mathcal O]=4-[g]<3$ containing three fields. Since each
field contributes at
least one unit of mass dimension, such an operator can only be built by inserting negative powers of
derivatives, and is therefore non-local. Demanding that the three-particle amplitude be finite on
real kinematics is thus equivalent to demanding that the underlying theory be local.

The three-particle amplitudes obtained here are the seeds of the whole tree-level $S$-matrix. Their
mutual consistency is itself constraining: demanding that the four-particle amplitude built from
them through the recursion relations of Section~\ref{sec:BCFW} be independent of the deformation
employed --- the so-called four-particle test~\cite{Benincasa:2007xk} --- forces the spin-$1$
couplings to satisfy the Jacobi identity, hence to be those of a Lie group, requires all spin-$2$
couplings to be equal, thereby deriving the equivalence principle rather than postulating it, and
excludes interacting massless particles of spin greater than $2$.

\subsubsection{Examples}

It is instructive to see the machinery at work in a few explicit cases. In each of them the
helicities alone fix the amplitude, and the mass dimension of the coupling follows without any
reference to a Lagrangian.

\paragraph{Cubic scalar coupling.}
With $(h_1,h_2,h_3)=(0,0,0)$ we have $h=0$ and $a=b=c=0$, so the
    amplitude is a constant,
    \begin{equation}
        \mathcal M (1^{0},2^{0},3^{0}) = \mu\,,\qquad [\mu]=1\,,
    \end{equation}
    the marginal case discussed above.

\paragraph{Yukawa coupling.}
With $(h_1,h_2,h_3)=(-\tfrac12,-\tfrac12,0)$ we have $h=-1$, the
configuration is holomorphic, and $(a,b,c)=(1,0,0)$. Thus,
    \begin{equation}
        \mathcal M (1^{-1/2},2^{-1/2},3^{0}) = y\,\agl{1}{2}\,,\qquad [y]=0\,,
    \end{equation}
correctly reproducing a dimensionless Yukawa coupling. The conjugate configuration
$(h_1,h_2,h_3)=(+\tfrac12,+\tfrac12,0)$ has $h=+1$ and gives $\overline y\,\sqr{1}{2}$.

\paragraph{Gauge boson coupling to matter.}
For a fermion,
an anti-fermion, and a gauge boson with $(h_1,h_2,h_3)=(-\tfrac12,+\tfrac12,-1)$, we have $h=-1$ and $(a,b,c)=(-1,0,2)$, so that
    \begin{equation}
        \mathcal M (1^{-1/2},2^{+1/2},3^{-1}) = g\,\frac{\agl{3}{1}^2}{\agl{1}{2}}\,,\qquad [g]=0\,,
    \end{equation}
which is the familiar three-point vertex of QED and QCD.

\paragraph{Yang--Mills.}
With $(h_1,h_2,h_3)=(-1,-1,+1)$ we have $h=-1$ and $(a,b,c)=(3,-1,-1)$. Thus,
    \begin{equation}
        \mathcal M (1^{-1},2^{-1},3^{+1}) = g\,\frac{\agl{1}{2}^3}{\agl{2}{3}\agl{3}{1}}\,,
        \qquad [g]=0\,,
    \end{equation}
which is the well-known three-gluon amplitude, whose color structure is left undetermined at this stage.

\paragraph{Gravity.}
With $(h_1,h_2,h_3)=(-2,-2,+2)$ we have $h=-2$ and $(a,b,c)=(6,-2,-2)$. Thus,
    \begin{equation}
        \mathcal M (1^{-2},2^{-2},3^{+2}) = \kappa\,\frac{\agl{1}{2}^6}{\agl{2}{3}^2\agl{3}{1}^2}\,,
        \qquad [\kappa]=-1\,,
    \end{equation}
the coupling being dimensionful and of order $1/M_{\text{P}}$. Note that this
amplitude is the square of the Yang--Mills one divided by the scalar one, an early
manifestation of the double-copy structure of gravitational amplitudes~\cite{Kawai:1985xq,Bern:2008qj,Bern:2010ue,Carrasco:2015iwa,Bern:2022wqg,Bern:2019prr}.

A pattern emerges from the gauge-like configurations $(-s,-s,+s)$, for which $h=-s$ and the coupling
carries dimension $[g]=1-s$: the cubic scalar coupling at $s=0$, the dimensionless gauge coupling at
$s=1$, and the gravitational coupling at $s=2$. Extrapolating further would require couplings of
increasingly negative dimension, and no consistent interaction of massless particles with
$s>2$ survives the consistency conditions mentioned above.

\paragraph{Axion/dilaton coupling to gauge bosons.}
With $(h_1,h_2,h_3)=(-1,-1,0)$ we have $h=-2$ and
$(a,b,c)=(2,0,0)$. Thus,
    \begin{equation}
        \mathcal M (1^{-1},2^{-1},3^{0}) = C\,\agl{1}{2}^2\,,\qquad [C]=-1\,,
    \end{equation}
which is the structure generated by the dimension-five operators 
$\phi F_{\mu\nu}F^{\mu\nu}$ and $\phi F_{\mu\nu}\tildee F^{\mu\nu}$.

\subsection{Recursion Relations for Higher-Point Amplitudes\label{sec:BCFW}}

The properties collected above are strong enough to determine tree-level amplitudes without any
reference to Feynman diagrams. Locality confines the singularities of an amplitude to the
factorization poles of Eq.~\eqref{eq:factorization}, and analyticity allows those poles to be
reached by deforming the external kinematics into the complex plane. Two meromorphic functions with the same poles and residues differ by a function that is regular
everywhere in the complex plane, and such a difference is absent when the deformed amplitude
vanishes at large complex deformation. An amplitude can therefore be reconstructed from on-shell
data of lower multiplicity whenever the deformation can be arranged to fall off in that limit. This is the content of the Britto--Cachazo--Feng--Witten (BCFW) recursion
relation~\cite{Britto:2004ap,Britto:2005fq} and of its generalizations~\cite{Risager:2005vk,Cohen:2010mi,Kampf:2012fn,Cheung:2014dqa,Cheung:2015ota,Cheung:2015cba}.

The construction starts from a complex deformation of two external legs. Choosing the pair $(i,j)$
and a complex parameter $z$, the shift
\begin{equation}\label{eq:BCFWshift}
    \lambda_i \longrightarrow \widehat\lambda_i(z) = \lambda_i - z\,\lambda_j\,, \qquad
    \tildee\lambda_j \longrightarrow \widehat{\tildee\lambda}_j(z) = \tildee\lambda_j + z\,\tildee\lambda_i\,,
\end{equation}
leaves $\tildee\lambda_i$ and $\lambda_j$ untouched and deforms the two momenta into
\begin{equation}
    \widehat p_i^{\dot\alpha\alpha}(z) = \tildee\lambda_i^{\dot\alpha}\big(\lambda_i - z\,\lambda_j\big)^{\alpha}\,,
    \qquad
    \widehat p_j^{\dot\alpha\alpha}(z) = \big(\tildee\lambda_j + z\,\tildee\lambda_i\big)^{\dot\alpha}\lambda_j^{\alpha}\,.
\end{equation}
Both stay light-like for every $z$, since each of them remains a product of two spinors, and their sum is independent of $z$. The deformed configuration therefore satisfies momentum conservation and the
on-shell conditions. The shifted momenta are complex, and this is what makes the three-particle
amplitudes of the previous subsection non-vanishing on the deformed kinematics, so that they can
serve as the seeds of the construction.

Let $\mathcal M_n(z)$ denote the amplitude evaluated on the deformed kinematics. It is a meromorphic
function of $z$ with simple poles wherever a channel momentum goes on shell. For a channel built out
of a subset $I$ of external legs, the shifted momentum $\widehat P_I(z)$ depends on $z$ only when $I$
contains exactly one of the two shifted legs, in which case $\widehat P_I^2(z)-m^2$ is linear in $z$ and
vanishes at a single point $z_I$. Applying Cauchy's theorem to $\mathcal M_n(z)/z$ on a small
contour around the origin and deforming it to infinity, the poles at $z_I$ are picked up with
opposite orientation and their residues are evaluated with Eq.~\eqref{eq:factorization}, giving
\begin{equation}\label{eq:recursion}
    \mathcal M_n = \sum_{I}\sum_{h}\,
    \mathcal M_{L}\big(z_I\big)\,\frac{1}{P_I^{2}-m^2}\,\mathcal M_{R}\big(z_I\big) + B\,,
\end{equation}
where the two subamplitudes are evaluated on the kinematics shifted to $z=z_I$, the sum over $h$
runs over the states of the exchanged particle, and $B$ is the contribution of the contour at
infinity. For the massless theories considered here the propagator reduces to $1/P_I^2$.

The recursion relation is Eq.~\eqref{eq:recursion} with $B$ dropped, and it is constructive, since
every amplitude on the right-hand side has fewer legs than $\mathcal M_n$. Discarding $B$ requires
$\mathcal M_n(z)\to0$ as $z\to\infty$, which is a property of the theory and of the helicities of the two
shifted legs. Admissible choices exist for gluon and graviton amplitudes, while for effective field
theories with derivative interactions the falloff is not guaranteed and has to be examined case by
case. Two further features are worth recording. Different admissible shifts reconstruct the same
amplitude through different sets of channels, so the resulting expressions look different and their
equivalence is not always manifest. Moreover, the individual terms of Eq.~\eqref{eq:recursion}
contain spurious poles, absent from the total, so locality is obscured in the final expression even
though it was the input of the derivation.

\subsubsection{Example: Parke--Taylor Formula}

For gluons the color and kinematic degrees of freedom can be disentangled. At tree level the
$n$-gluon amplitude admits the color decomposition~\cite{Mangano:1990by,Berends:1987cv}
\begin{equation}\label{eq:colordec}
    \mathcal M_n\big(\{p_i,h_i,a_i\}\big) = g^{n-2}\sum_{\sigma\in S_{n-1}}
    \Tr(T^{a_1}T^{a_{\sigma(2)}}\cdots T^{a_{\sigma(n)}})\,
    A_n\big(1^{h_1},\sigma(2)^{h_{\sigma(2)}},\dotsc,\sigma(n)^{h_{\sigma(n)}}\big)\,,
\end{equation}
where $T^a$ are the generators of the gauge group in the fundamental representation, normalized as
$\Tr T^aT^b =\delta^{ab}$, and the sum runs over the permutations of the labels $2,\dotsc,n$. The functions $A_n$ are the \textit{color-ordered} or \textit{partial amplitudes}. They are gauge
invariant, invariant under cyclic relabeling of their arguments, and singular only in channels built
out of cyclically adjacent momenta. This last property is what makes the shift of two adjacent legs
natural, since the recursion then receives contributions from $n-3$ channels only.

Tree-level partial amplitudes with all helicities equal, or with a single helicity differing from
the others, vanish for $n\ge4$. The first non-vanishing configurations carry two negative helicities
and are called \textit{maximally-helicity-violating} (MHV). For them the recursion collapses to the
remarkably compact Parke--Taylor formula~\cite{Parke:1986gb}
\begin{equation}\label{eq:parketaylor}
    A_n\big(1^+,\dotsc,i^-,\dotsc,j^-,\dotsc,n^+\big)
    = \frac{\agl{i}{j}^4}{\agl{1}{2}\agl{2}{3}\cdots\agl{n-1\,}{n}\agl{n}{1}}\,,
\end{equation}
with the conjugate anti-MHV configuration obtained by exchanging angle and square brackets.
Eq.~\eqref{eq:parketaylor} passes the consistency checks established above. Its little-group
weight on the leg $i$ is $t_i^{4-2}=t_i^{2}$, corresponding to $h_i=-1$, while a positive-helicity
leg $k$ appears only in two denominator factors and carries $t_k^{-2}$. Its mass dimension is
$4-n$, in agreement with Eq.~\eqref{eq:massdim} for a dimensionless coupling. At $n=3$ it reduces to
$\agl{1}{2}^3/(\agl{2}{3}\agl{3}{1})$, reproducing the three-gluon amplitude obtained from the
bootstrap of the previous subsection. A single expression, built recursively out of three-point
amplitudes that were themselves fixed by helicity weights and dimensional analysis, replaces a
number of Feynman diagrams that grows factorially with $n$.
